\documentclass[aspectratio=169,10pt]{beamer}

% Shared Beamer setup for Fall 2026 course slides.
\mode<presentation>{
  \usetheme{Madrid}
  \usecolortheme{default}
}

\definecolor{CalPolyGreen}{HTML}{154734}
\definecolor{CalPolyGold}{HTML}{BD8B13}
\definecolor{DeckInk}{HTML}{1F2933}
\definecolor{DeckMuted}{HTML}{5F6B7A}

\setbeamercolor{structure}{fg=CalPolyGreen}
\setbeamercolor{palette primary}{bg=CalPolyGreen,fg=white}
\setbeamercolor{palette secondary}{bg=DeckInk,fg=white}
\setbeamercolor{palette tertiary}{bg=CalPolyGold,fg=DeckInk}
\setbeamercolor{frametitle}{bg=CalPolyGreen,fg=white}
\setbeamercolor{title}{fg=CalPolyGreen}
\setbeamercolor{normal text}{fg=DeckInk,bg=white}
\setbeamercolor{block title}{bg=CalPolyGreen,fg=white}
\setbeamercolor{block body}{bg=black!3,fg=DeckInk}

\setbeamertemplate{navigation symbols}{}
\setbeamertemplate{footline}[frame number]
\setbeamertemplate{itemize item}{\color{CalPolyGold}$\blacktriangleright$}
\setbeamertemplate{itemize subitem}{\color{CalPolyGreen}$\bullet$}

\newcommand{\CourseNumber}{Course}
\newcommand{\CourseTitle}{Course Title}
\newcommand{\LectureNumber}{0}
\newcommand{\LectureTitle}{Lecture Title}
\newcommand{\LectureDate}{Date}
\newcommand{\CourseUnit}{Unit}

\newcommand{\coursenumber}[1]{\renewcommand{\CourseNumber}{#1}}
\newcommand{\coursetitle}[1]{\renewcommand{\CourseTitle}{#1}}
\newcommand{\lecturenumber}[1]{\renewcommand{\LectureNumber}{#1}}
\newcommand{\lecturetitle}[1]{\renewcommand{\LectureTitle}{#1}}
\newcommand{\lecturedate}[1]{\renewcommand{\LectureDate}{#1}}
\newcommand{\courseunit}[1]{\renewcommand{\CourseUnit}{#1}}

\author{Kirk Duran}
\institute{California Polytechnic State University, San Luis Obispo}
\date{\LectureDate}

\newcommand{\makedecktitle}{
  \begin{frame}[plain]
    \vspace{0.25in}
    {\usebeamerfont{title}\color{CalPolyGreen}\LectureTitle\par}
    \vspace{0.18in}
    {\large \CourseNumber{}: \CourseTitle\par}
    \vspace{0.08in}
    {\large Lecture \LectureNumber{} -- \LectureDate\par}
    \vfill
    {\small Kirk Duran \textbar{} Cal Poly, San Luis Obispo\par}
  \end{frame}
}

\newcommand{\placeholdernote}{
  \begin{block}{Placeholder}
    Replace these bullets with the final lecture material, examples, and in-class activities.
  \end{block}
}

\AtBeginSection[]{
  \begin{frame}{Roadmap}
    \tableofcontents[currentsection]
  \end{frame}
}

\usepackage{tikz}

\graphicspath{{../assets/lecture-05/}}

% If an image file is not yet available, skip the visual slot.
\newcommand{\lectureimage}[2][1.75in]{%
  \IfFileExists{../assets/lecture-05/#2}{%
    \includegraphics[width=\linewidth,height=#1,keepaspectratio]{#2}%
  }{%
  }%
}

\setbeamersize{text margin left=0.42in,text margin right=0.42in}

\coursenumber{CS 4880}
\coursetitle{Artificial Intelligence}
\lecturenumber{5}
\lecturetitle{Uninformed Search II: Action Costs and Uniform-Cost Search}
\lecturedate{September 2, 2026}
\courseunit{Search and Problem Solving}

\begin{document}

\makedecktitle

\begin{frame}[shrink=20]{Today}
  \begin{itemize}
    \item How Micromouse exposed limitations in a seemingly solved search formulation.
    \item Why richer action sets and physical motion constraints change the state space.
    \item How kinematics, velocity, and environmental dynamics affect feasible actions.
    \item Why unequal actions require an explicit action-cost model.
    \item How Dijkstra's algorithm / Uniform-Cost Search finds minimum-cost solutions.
    \item Where minimum-cost graph search appears as a baseline or subroutine in modern AI systems.
  \end{itemize}

  \begin{keyidea}
    Search quality depends not only on the search algorithm, but also on how we define the states, actions, transitions, and objective.
  \end{keyidea}
\end{frame}

\sectionframe{Part 1: Micromouse Revisited --- When the Model Is Too Simple}

\begin{frame}[shrink=20]{Last Time: A Deliberately Simple Micromouse Model}
  \begin{columns}[T,onlytextwidth]
    \begin{column}{0.54\textwidth}
      \begin{block}{Discrete search abstraction}
        \begin{itemize}
          \item State: current maze cell
          \item Actions: move to an adjacent open cell
          \item Goal: reach a center cell
          \item Objective: minimize the number of actions
        \end{itemize}
      \end{block}

      \begin{itemize}
        \item BFS could find a minimum-depth route once the maze was known.
        \item Flood-fill style methods exploited the same grid structure efficiently.
      \end{itemize}
    \end{column}

    \begin{column}{0.40\textwidth}
      % Search direction: top-down Micromouse maze rendered as an orthogonal grid with a highlighted center goal.
      \lectureimage[2.45in]{slide4.png}
    \end{column}
  \end{columns}

  \begin{keyidea}
    The abstraction was useful because it removed physical details. Today we ask what happens when some of those details become important.
  \end{keyidea}
\end{frame}

\begin{frame}[shrink=20]{A Search Algorithm Can Only Use the Actions We Define}
  \begin{columns}[T,onlytextwidth]
    \begin{column}{0.48\textwidth}
      \begin{block}{Early orthogonal action model}
        \centering
        \texttt{FORWARD}\\
        \texttt{TURN LEFT 90$^\circ$}\\
        \texttt{TURN RIGHT 90$^\circ$}
      \end{block}

      \begin{itemize}
        \item Paths are constructed from horizontal and vertical segments.
        \item A change in direction is represented as a 90-degree turn.
        \item Diagonal trajectories are not expressible in the model.
      \end{itemize}
    \end{column}

    \begin{column}{0.46\textwidth}
      % Search direction: side-by-side maze routes, one orthogonal stair-step route and one unavailable diagonal shortcut.
      \lectureimage[2.55in]{slide5.png}
    \end{column}
  \end{columns}

  \begin{keyidea}
    A solution cannot be discovered if the action representation makes that solution impossible to express.
  \end{keyidea}
\end{frame}

\begin{frame}[shrink=20]{Micromouse Risked Becoming a ``Solved'' Problem}
  \begin{columns}[T,onlytextwidth]
    \begin{column}{0.56\textwidth}
      \begin{itemize}
        \item By the late 1980s, reliable maze discovery and shortest-path methods were well understood.
        \item The research community began to question whether the competition still presented substantial new technical challenges.
        \item The apparent problem had become: discover the maze, compute a shortest orthogonal route, then execute it faster.
      \end{itemize}

      \begin{block}{But this conclusion depended on the formulation}
        If states and actions remain fixed, performance eventually becomes dominated by incremental implementation improvements.
      \end{block}
    \end{column}

    \begin{column}{0.38\textwidth}
      % Search direction: historical Micromouse competition timeline showing apparent stagnation in the late 1980s.
      \lectureimage[2.48in]{slide6.jpeg}
    \end{column}
  \end{columns}

  \note{[Sources] Veritasium, ``The Fastest Maze-Solving Competition On Earth,'' https://www.youtube.com/watch?v=ZMQbHMgK2rw; Peter Harrison, Micromouse Online history, https://micromouseonline.com/micromouse-book/history/.}
\end{frame}

\begin{frame}[shrink=20]{Mitee Mouse III Changed the Action Space}
  \begin{columns}[T,onlytextwidth]
    \begin{column}{0.54\textwidth}
      \begin{itemize}
        \item Mitee Mouse III demonstrated diagonal running rather than restricting every route to orthogonal grid segments.
        \item A 45-degree transition could enter or leave a diagonal corridor through the maze.
        \item The maze did not change; the \textbf{agent's action set} changed.
      \end{itemize}

      \begin{block}{New possibilities}
        \centering
        orthogonal: $\rightarrow\;\downarrow\;\rightarrow$\\[0.10in]
        diagonal: $\rightarrow\;\searrow\;\searrow\;\rightarrow$
      \end{block}
    \end{column}

    \begin{column}{0.40\textwidth}
      % Search direction: Mitee Mouse III or a Micromouse route illustrating 45-degree diagonal movement through square maze cells.
      \lectureimage[2.48in]{slide7.png}
    \end{column}
  \end{columns}

  \begin{keyidea}
    Expanding \texttt{ACTIONS(s)} can reveal qualitatively new solutions, not merely faster implementations of old ones.
  \end{keyidea}

  \note{[Sources] Veritasium, ``The Fastest Maze-Solving Competition On Earth,'' https://www.youtube.com/watch?v=ZMQbHMgK2rw; Micromouse Online, diagonal path conventions and turn primitives, https://micromouseonline.com/2014/05/05/diagonal-paths-micromouse-using-state-machine/.}
\end{frame}

\begin{frame}[shrink=20]{Richer Actions Create a Richer Search Space}
  \begin{columns}[T,onlytextwidth]
    \begin{column}{0.46\textwidth}
      \begin{block}{Orthogonal model}
        \begin{itemize}
          \item forward
          \item 90-degree left
          \item 90-degree right
        \end{itemize}
      \end{block}
    \end{column}

    \begin{column}{0.50\textwidth}
      \begin{block}{Richer motion primitives}
        \begin{itemize}
          \item straight motion
          \item 90-degree smooth turns
          \item enter diagonal at 45 degrees
          \item diagonal travel
          \item leave diagonal at 45 degrees
          \item larger composite turns
        \end{itemize}
      \end{block}
    \end{column}
  \end{columns}

  \begin{itemize}
    \item More actions increase the branching factor.
    \item More actions may also reduce solution length or execution time dramatically.
    \item Search complexity and representational power grow together.
  \end{itemize}

  \begin{keyidea}
    The action model is part of the intelligence: it determines what futures the agent is capable of considering.
  \end{keyidea}
\end{frame}

\begin{frame}[shrink=20]{Orientation Now Becomes Part of the State}
  \begin{columns}[T,onlytextwidth]
    \begin{column}{0.53\textwidth}
      Two mice can occupy the same maze cell but face different directions.

      \begin{block}{Too little information}
        \centering
        $s=(r,c)$
      \end{block}

      \begin{block}{Better state representation}
        \centering
        $s=(r,c,\theta)$
      \end{block}

      \begin{itemize}
        \item $r,c$ describe location.
        \item $\theta$ describes orientation or heading.
        \item Available next actions depend on all three.
      \end{itemize}
    \end{column}

    \begin{column}{0.41\textwidth}
    \end{column}
  \end{columns}

  \begin{keyidea}
    A state must contain enough information for the transition model to predict the consequences of future actions.
  \end{keyidea}
\end{frame}

\sectionframe{Part 2: Motion Constraints --- From Grid Search to Physical Planning}

\begin{frame}[shrink=20]{A Car Cannot Simply ``Move Left''}
  \begin{columns}[T,onlytextwidth]
    \begin{column}{0.54\textwidth}
      A grid model might allow an action such as:

      \begin{block}{Abstract action}
        \centering
        \texttt{LEFT}
      \end{block}

      A real car cannot translate sideways. Steering changes the vehicle's heading while it moves forward, producing an arc.

      \begin{itemize}
        \item Position changes continuously.
        \item Heading changes continuously.
        \item The vehicle has a nonzero turning radius.
      \end{itemize}
    \end{column}

    \begin{column}{0.40\textwidth}
    \end{column}
  \end{columns}

  \begin{keyidea}
    Path planning must respect the motions the agent can physically execute, not merely geometric adjacency.
  \end{keyidea}
\end{frame}

\begin{frame}[shrink=20]{Kinematic Constraints Restrict Reachable States}
  \begin{cpdefinition}
    \textbf{Kinematics} describes the geometry of motion---how position and orientation evolve---without modeling the forces that cause the motion.
  \end{cpdefinition}

  \begin{columns}[T,onlytextwidth]
    \begin{column}{0.50\textwidth}
      A simple vehicle state may include:
      \[
        s=(x,y,\theta)
      \]

      where:
      \begin{itemize}
        \item $x,y$ are position;
        \item $\theta$ is heading.
      \end{itemize}

      Steering limits may impose a minimum turning radius.
    \end{column}

    \begin{column}{0.44\textwidth}
    \end{column}
  \end{columns}
\end{frame}

\begin{frame}[shrink=20]{A Minimal Kinematic Bicycle Model}
  \begin{columns}[T,onlytextwidth]
    \begin{column}{0.50\textwidth}
      A common simplified model is:
      \[
        \dot{x}=v\cos\theta
      \]
      \[
        \dot{y}=v\sin\theta
      \]
      \[
        \dot{\theta}=\frac{v}{L}\tan\delta
      \]

      where:
      \begin{itemize}
        \item $v$ is forward velocity,
        \item $L$ is wheelbase,
        \item $\delta$ is steering angle.
      \end{itemize}
    \end{column}

    \begin{column}{0.44\textwidth}
      \begin{block}{Interpretation}
        \begin{itemize}
          \item Steering does not instantaneously rotate the car.
          \item Turning rate depends on both speed and steering angle.
          \item Some geometrically short paths are not kinematically feasible.
        \end{itemize}
      \end{block}
    \end{column}
  \end{columns}

  \begin{keyidea}
    The transition model must encode the system's motion constraints; otherwise search can return physically impossible plans.
  \end{keyidea}
\end{frame}

\begin{frame}[shrink=20]{Velocity Can Change Which Actions Are Feasible}
  \begin{columns}[T,onlytextwidth]
    \begin{column}{0.54\textwidth}
      Consider two cars at the same position and heading:

      \begin{center}
      \renewcommand{\arraystretch}{1.25}
      \begin{tabular}{c|c}
        \textbf{Vehicle} & \textbf{Velocity} \\\hline
        A & 5 mph \\
        B & 60 mph
      \end{tabular}
      \end{center}

      A maneuver that is feasible for A may be impossible or unsafe for B.

      \begin{block}{Richer state}
        \centering
        $s=(x,y,\theta,v)$
      \end{block}
    \end{column}

    \begin{column}{0.40\textwidth}
    \end{column}
  \end{columns}

  \begin{keyidea}
    If future actions depend on velocity, velocity belongs in the state representation.
  \end{keyidea}
\end{frame}

\begin{frame}[shrink=20]{Stopping Distance Grows Rapidly with Speed}
  Under an idealized constant-deceleration model:
  \[
    d_{\text{stop}} \approx \frac{v^2}{2a}
  \]

  \begin{columns}[T,onlytextwidth]
    \begin{column}{0.47\textwidth}
      \begin{block}{Low speed}
        At 5 mph, the kinetic energy and required braking distance are comparatively small.
      \end{block}
    \end{column}

    \begin{column}{0.47\textwidth}
      \begin{block}{High speed}
        At 60 mph, the same braking capability cannot stop the vehicle within the same short distance.
      \end{block}
    \end{column}
  \end{columns}

  \begin{keyidea}
    Search over physical systems may require state-dependent actions: \texttt{ACTIONS(s)} changes as speed and configuration change.
  \end{keyidea}
\end{frame}

\begin{frame}[shrink=20]{Beyond Kinematics: Dynamics and External Conditions}
  \begin{cpdefinition}
    \textbf{Dynamics} concerns how forces, mass, inertia, friction, and actuation determine motion.
  \end{cpdefinition}

  \begin{columns}[T,onlytextwidth]
    \begin{column}{0.50\textwidth}
      The same steering command can produce different outcomes under different conditions:
      \begin{itemize}
        \item dry pavement,
        \item wet pavement,
        \item ice,
        \item loose gravel.
      \end{itemize}

      A simple friction-limited model gives roughly:
      \[
        a_{\max}\approx \mu g
      \]
    \end{column}

    \begin{column}{0.44\textwidth}
    \end{column}
  \end{columns}
\end{frame}

\begin{frame}[shrink=20]{The Transition Model Depends on the Physical World}
  Previously:
  \[
    \mathrm{RESULT}(s,a)=s'
  \]

  looked almost like a table lookup.

  For a physical agent, the result may depend on:
  \begin{itemize}
    \item current position and orientation;
    \item current velocity;
    \item steering and acceleration limits;
    \item friction and traction;
    \item environmental conditions.
  \end{itemize}

  \begin{keyidea}
    Problem formulation is a modeling decision: include enough physical detail to predict useful consequences without making planning computationally intractable.
  \end{keyidea}
\end{frame}

\begin{frame}[shrink=20]{Micromouse Faced the Same Physical Limits}
  \begin{columns}[T,onlytextwidth]
    \begin{column}{0.55\textwidth}
      As Micromouse speeds increased:
      \begin{itemize}
        \item turning required greater lateral acceleration;
        \item tire traction became a limiting factor;
        \item wheel cleanliness and surface friction mattered;
        \item later designs introduced suction fans to increase normal force and cornering capability.
      \end{itemize}

      \begin{block}{The software and hardware co-evolved}
        Better planning was useful only when the robot could physically execute the resulting trajectory.
      \end{block}
    \end{column}

    \begin{column}{0.39\textwidth}
    \end{column}
  \end{columns}

  \note{[Sources] Veritasium, ``The Fastest Maze-Solving Competition On Earth,'' https://www.youtube.com/watch?v=ZMQbHMgK2rw; Micromouse competition discussions summarized at https://hackaday.com/2023/05/28/the-fascinating-evolution-of-micromouse/.}
\end{frame}

\sectionframe{Part 3: From Feasibility to Preference --- Action Costs}

\begin{frame}[shrink=20]{Feasible Does Not Mean Equally Desirable}
  Suppose two routes both reach the goal.

  \begin{columns}[T,onlytextwidth]
    \begin{column}{0.47\textwidth}
      \begin{block}{Route A}
        \begin{itemize}
          \item 5 actions
          \item many sharp turns
          \item repeated deceleration
          \item execution time: 8.2 s
        \end{itemize}
      \end{block}
    \end{column}

    \begin{column}{0.47\textwidth}
      \begin{block}{Route B}
        \begin{itemize}
          \item 7 actions
          \item mostly straight motion
          \item higher sustained speed
          \item execution time: 6.4 s
        \end{itemize}
      \end{block}
    \end{column}
  \end{columns}

  \begin{keyidea}
    Minimizing the number of actions is only appropriate when all actions are effectively equivalent.
  \end{keyidea}
\end{frame}

\begin{frame}[shrink=20]{Shortest Path Is Not Necessarily Fastest Path}
  \begin{columns}[T,onlytextwidth]
    \begin{column}{0.54\textwidth}
      The Micromouse competition provides a concrete example:
      \begin{itemize}
        \item a geometrically longer route can contain fewer or gentler turns;
        \item fewer costly turns can permit a higher sustained velocity;
        \item therefore a longer route can have lower execution time.
      \end{itemize}

      \begin{block}{Objective matters}
        \centering
        minimum distance $\neq$ minimum time
      \end{block}
    \end{column}

    \begin{column}{0.40\textwidth}
      % Search direction: two Micromouse routes through the same maze, shorter route with many turns versus longer route with fewer turns and lower time.
      \lectureimage[2.50in]{slide21.png}
    \end{column}
  \end{columns}

  \note{[Sources] Veritasium, ``The Fastest Maze-Solving Competition On Earth,'' https://www.youtube.com/watch?v=ZMQbHMgK2rw.}
\end{frame}

\begin{frame}[shrink=20]{Action Cost}
  \begin{cpdefinition}
    The \textbf{action cost} (or \textbf{step cost}) assigns a numerical cost to performing action $a$ in state $s$ and reaching state $s'$.
  \end{cpdefinition}

  \[
    c(s,a,s')
  \]

  Action cost can represent different objectives:
  \begin{itemize}
    \item execution time;
    \item physical distance;
    \item energy consumption;
    \item monetary cost;
    \item risk or exposure;
    \item a designed combination of several factors.
  \end{itemize}

  \begin{keyidea}
    The cost function defines what the search algorithm means by a ``better'' solution.
  \end{keyidea}
\end{frame}

\begin{frame}[shrink=20]{Path Cost Accumulates Action Costs}
  For a path
  \[
    s_0 \xrightarrow{a_1} s_1 \xrightarrow{a_2} s_2 \cdots \xrightarrow{a_n} s_n,
  \]

  define cumulative path cost:
  \[
    g(n)=\sum_{i=1}^{n} c(s_{i-1},a_i,s_i).
  \]

  \begin{columns}[T,onlytextwidth]
    \begin{column}{0.47\textwidth}
      \begin{block}{Lecture 5 objective}
        Minimize solution depth:
        \[
          d(n)
        \]
      \end{block}
    \end{column}

    \begin{column}{0.47\textwidth}
      \begin{block}{Lecture 6 objective}
        Minimize cumulative path cost:
        \[
          g(n)
        \]
      \end{block}
    \end{column}
  \end{columns}
\end{frame}

\begin{frame}[shrink=20]{Depth and Cost Can Disagree}
  \begin{columns}[T,onlytextwidth]
    \begin{column}{0.54\textwidth}
      Consider two solutions:

      \begin{block}{Shallow solution}
        \centering
        $S\rightarrow A\rightarrow G$\\
        depth $=2$, cost $=16$
      \end{block}

      \begin{block}{Deeper solution}
        \centering
        $S\rightarrow B\rightarrow C\rightarrow G$\\
        depth $=3$, cost $=6$
      \end{block}
    \end{column}

    \begin{column}{0.40\textwidth}
    \end{column}
  \end{columns}

  \begin{keyidea}
    Once edge costs differ, the shallowest goal may not be the best goal.
  \end{keyidea}
\end{frame}

\begin{frame}[shrink=20]{Why Breadth-First Search Is No Longer Enough}
  BFS prioritizes:
  \[
    \text{smallest depth}
  \]

  but weighted search requires:
  \[
    \text{smallest accumulated cost}
  \]

  \begin{center}
  \renewcommand{\arraystretch}{1.30}
  \begin{tabular}{c|c|c}
    \textbf{Frontier node} & \textbf{Depth} & \textbf{Path cost $g$} \\\hline
    A & 2 & 3 \\
    B & 2 & 50 \\
    C & 3 & 4
  \end{tabular}
  \end{center}

  BFS sees A and B as equally shallow. A cost-aware search should prefer A, then possibly C, before B.

  \begin{keyidea}
    We need a frontier policy ordered by cumulative path cost rather than depth.
  \end{keyidea}
\end{frame}

\sectionframe{Part 4: Dijkstra's Algorithm and Uniform-Cost Search}

\begin{frame}[shrink=20]{Edsger W. Dijkstra}
  \begin{columns}[T,onlytextwidth]
    \begin{column}{0.55\textwidth}
      \begin{itemize}
        \item Dijkstra approached computing as a mathematical discipline centered on precise algorithms and proofs.
        \item In 1959, he published \emph{A Note on Two Problems in Connexion with Graphs}.
        \item One of the algorithms in that paper computes shortest paths in a graph with non-negative edge weights.
      \end{itemize}

      \begin{block}{Central idea}
        Repeatedly select the reachable unsettled vertex with the smallest known cumulative distance.
      \end{block}
    \end{column}

    \begin{column}{0.39\textwidth}
      % Search direction: portrait of Edsger Dijkstra paired with the 1959 paper title or weighted graph sketch.
      \lectureimage[2.50in]{slide27.jpg}
    \end{column}
  \end{columns}

  \note{[Sources] E. W. Dijkstra, ``A Note on Two Problems in Connexion with Graphs,'' Numerische Mathematik 1, 1959, pp. 269--271.}
\end{frame}

\begin{frame}[shrink=20]{Dijkstra vs. Uniform-Cost Search}
  \begin{columns}[T,onlytextwidth]
    \begin{column}{0.47\textwidth}
      \begin{block}{Dijkstra's algorithm}
        Typically presented as a single-source shortest-path algorithm that computes minimum costs from the source to reachable vertices.
      \end{block}
    \end{column}

    \begin{column}{0.47\textwidth}
      \begin{block}{Uniform-Cost Search (UCS)}
        The AI search formulation: repeatedly expand the frontier node with minimum path cost and stop when a goal is removed for expansion.
      \end{block}
    \end{column}
  \end{columns}

  Both use the same ordering principle:
  \[
    \text{priority}(n)=g(n).
  \]

  \begin{keyidea}
    UCS is the goal-directed search form of Dijkstra's minimum-cost principle.
  \end{keyidea}
\end{frame}

\begin{frame}[shrink=20]{Frontier Policy: Always Expand the Cheapest Path So Far}
  \begin{block}{Uniform-Cost rule}
    \centering
    remove the frontier node with minimum $g(n)$
  \end{block}

  \begin{columns}[T,onlytextwidth]
    \begin{column}{0.30\textwidth}
      \begin{block}{BFS}
        FIFO queue\\[0.06in]
        priority $=$ depth
      \end{block}
    \end{column}

    \begin{column}{0.30\textwidth}
      \begin{block}{DFS}
        LIFO stack\\[0.06in]
        priority $=$ depth-first order
      \end{block}
    \end{column}

    \begin{column}{0.34\textwidth}
      \begin{block}{UCS / Dijkstra}
        min-priority queue\\[0.06in]
        priority $=g(n)$
      \end{block}
    \end{column}
  \end{columns}

  \begin{keyidea}
    The algorithm remains uninformed: $g(n)$ measures only the cost already incurred, not how close $n$ appears to the goal.
  \end{keyidea}
\end{frame}

\begin{frame}[shrink=20]{Worked Example: Weighted Search Graph}
  \begin{columns}[T,onlytextwidth]
    \begin{column}{0.46\textwidth}
      We will use the same graph for the next several slides.

      \begin{block}{Relevant edge costs}
        \begin{itemize}
          \item $S\!\rightarrow\!A:2$
          \item $S\!\rightarrow\!B:5$
          \item $A\!\rightarrow\!C:2$
          \item $A\!\rightarrow\!D:5$
          \item $B\!\rightarrow\!D:1$
          \item $C\!\rightarrow\!D:1$
          \item $C\!\rightarrow\!G:7$
          \item $D\!\rightarrow\!G:2$
        \end{itemize}
      \end{block}
    \end{column}

    \begin{column}{0.48\textwidth}
      \centering
      \begin{tikzpicture}[
        >=stealth,
        every node/.style={font=\small},
        state/.style={
          circle,
          draw=CalPolyGreen,
          very thick,
          fill=CalPolyBeige,
          minimum size=0.38in,
          inner sep=0pt
        },
        edge/.style={->, very thick, draw=DeckMuted},
        cost/.style={fill=white, inner sep=1.5pt, font=\scriptsize}
      ]
        \node[state] (S) at (0,1.1) {S};
        \node[state] (A) at (1.5,2.0) {A};
        \node[state] (B) at (1.5,0.2) {B};
        \node[state] (C) at (3.0,2.0) {C};
        \node[state] (D) at (3.0,0.2) {D};
        \node[state] (G) at (4.5,1.1) {G};

        \draw[edge] (S) -- node[cost, above left] {2} (A);
        \draw[edge] (S) -- node[cost, below left] {5} (B);
        \draw[edge] (A) -- node[cost, above] {2} (C);
        \draw[edge] (A) -- node[cost, right] {5} (D);
        \draw[edge] (B) -- node[cost, below] {1} (D);
        \draw[edge] (C) -- node[cost, right] {1} (D);
        \draw[edge] (C) -- node[cost, above right] {7} (G);
        \draw[edge] (D) -- node[cost, below right] {2} (G);
      \end{tikzpicture}
    \end{column}
  \end{columns}
\end{frame}

\begin{frame}[shrink=20]{Dijkstra Step 0: Initialize the Frontier}
  \begin{columns}[T,onlytextwidth]
    \begin{column}{0.50\textwidth}
      Initial values:
      \[
        g(S)=0
      \]
      \[
        g(v)=\infty \quad \text{for all other vertices}
      \]

      Frontier:
      \begin{block}{Priority queue}
        \centering
        $(S,0)$
      \end{block}

      Parent:
      \[
        \mathrm{parent}[S]=\mathrm{None}
      \]
    \end{column}

    \begin{column}{0.44\textwidth}
    \end{column}
  \end{columns}
\end{frame}

\begin{frame}[shrink=20]{Dijkstra Step 1: Expand \texttt{S}}
  Remove the minimum-cost frontier node:
  \[
    S:0
  \]

  Relax its outgoing edges:
  \[
    g(A)=0+2=2
  \]
  \[
    g(B)=0+5=5
  \]

  \begin{columns}[T,onlytextwidth]
    \begin{column}{0.47\textwidth}
      \begin{block}{Frontier}
        \centering
        $A:2$\\
        $B:5$
      \end{block}
    \end{column}

    \begin{column}{0.47\textwidth}
    \end{column}
  \end{columns}
\end{frame}

\begin{frame}[shrink=20]{Dijkstra Step 2: Expand \texttt{A}}
  The cheapest frontier node is:
  \[
    A:2
  \]

  Relax successors:
  \[
    g(C)=2+2=4
  \]
  \[
    g(D)=2+5=7
  \]

  New frontier:
  \begin{block}{Priority queue ordered by $g$}
    \centering
    $C:4,\quad B:5,\quad D:7$
  \end{block}

  \begin{keyidea}
    Search order is determined by cumulative cost, not by how visually close a node is to the goal.
  \end{keyidea}
\end{frame}

\begin{frame}[shrink=20]{Dijkstra Step 3: Relaxation Improves a Known Path}
  Expand:
  \[
    C:4
  \]

  From $C$:
  \[
    \text{candidate}(D)=4+1=5
  \]

  but our current best value is:
  \[
    g(D)=7.
  \]

  Since $5<7$:
  \[
    g(D)\leftarrow 5,
    \qquad
    \mathrm{parent}[D]\leftarrow C.
  \]

  \begin{keyidea}
    \textbf{Relaxation} replaces a known route when a newly discovered route reaches the same state more cheaply.
  \end{keyidea}
\end{frame}

\begin{frame}[shrink=20]{Dijkstra Step 4: Reach the Goal at Minimum Cost}
  After the remaining relaxations:
  \begin{itemize}
    \item $B$ has cost $5$;
    \item $D$ has cost $5$;
    \item expanding $D$ produces $G$ with cost $5+2=7$;
    \item the earlier route $C\rightarrow G$ would cost $4+7=11$.
  \end{itemize}

  When $G:7$ becomes the minimum item removed from the priority queue, the solution is final.

  \begin{block}{Minimum-cost path}
    \centering
    $S\rightarrow A\rightarrow C\rightarrow D\rightarrow G$\\[0.08in]
    total cost $=2+2+1+2=7$
  \end{block}

  \begin{keyidea}
    For non-negative edge costs, the goal is safe to return when it is removed as the minimum-cost frontier node---not merely when it is first generated.
  \end{keyidea}
\end{frame}

\begin{frame}[shrink=20]{Uniform-Cost Search Pseudocode}
  \begin{block}{Algorithm}
  \ttfamily\small
  frontier $\leftarrow$ PriorityQueue()\\
  frontier.insert(initial\_state, priority=0)\\
  cost\_so\_far[initial\_state] $\leftarrow$ 0\\
  came\_from[initial\_state] $\leftarrow$ None\\[0.06in]

  while frontier is not empty:\\
  \quad state $\leftarrow$ frontier.remove\_min()\\[0.03in]
  \quad if goal\_test(state):\\
  \qquad return reconstruct\_path(came\_from, state)\\[0.03in]
  \quad for action in ACTIONS(state):\\
  \qquad child $\leftarrow$ RESULT(state, action)\\
  \qquad new\_cost $\leftarrow$ cost\_so\_far[state] + ACTION\_COST(state, action, child)\\
  \qquad if child is new or new\_cost $<$ cost\_so\_far[child]:\\
  \qquad\quad cost\_so\_far[child] $\leftarrow$ new\_cost\\
  \qquad\quad came\_from[child] $\leftarrow$ state\\
  \qquad\quad frontier.insert(child, priority=new\_cost)\\[0.03in]
  return failure
  \end{block}
\end{frame}

\begin{frame}[shrink=20]{Why Is the Cheapest Frontier Node Safe to Finalize?}
  Suppose $n$ is the frontier node with minimum cost $g(n)$.

  Could an undiscovered path later reach $n$ for less?

  \begin{enumerate}
    \item Any alternate path must first pass through some not-yet-expanded frontier state $m$.
    \item By construction, $g(m)\ge g(n)$.
    \item All remaining edge costs are non-negative.
    \item Continuing from $m$ cannot reduce the accumulated cost below $g(m)$.
  \end{enumerate}

  Therefore no later path can improve $n$ once it is removed as the cheapest frontier state.

  \begin{keyidea}
    Non-negative edge costs are what make Dijkstra's greedy finalization step correct.
  \end{keyidea}
\end{frame}

\begin{frame}[shrink=20]{Properties of Dijkstra / Uniform-Cost Search}
  \begin{center}
  \renewcommand{\arraystretch}{1.32}
  \begin{tabular}{l|l}
    \textbf{Property} & \textbf{Result} \\\hline
    Complete & Yes on a finite graph with non-negative edge costs \\
    Optimal & Yes; returns a minimum-cost path \\
    Frontier & Min-priority queue ordered by $g(n)$ \\
    Time & $O((|V|+|E|)\log |V|)$ with adjacency list + binary heap \\
    Auxiliary search space & $O(|V|)$ for distances, parents, and frontier metadata
  \end{tabular}
  \end{center}

  \begin{alertblock}{Restriction}
    Dijkstra's finalization rule is not valid in the presence of negative-cost edges.
  \end{alertblock}

  \note{[Sources] E. W. Dijkstra, 1959; standard binary-heap implementation analysis as presented in algorithms texts such as Cormen et al., Introduction to Algorithms.}
\end{frame}

\begin{frame}[shrink=20]{What Dijkstra Pays for Its Guarantee}
  \begin{columns}[T,onlytextwidth]
    \begin{column}{0.47\textwidth}
      \begin{block}{Strengths}
        \begin{itemize}
          \item complete under standard assumptions;
          \item minimum-cost guarantee;
          \item domain-independent;
          \item handles arbitrary non-negative edge costs.
        \end{itemize}
      \end{block}
    \end{column}

    \begin{column}{0.47\textwidth}
      \begin{block}{Costs}
        \begin{itemize}
          \item potentially large frontier;
          \item stores best known costs and parents;
          \item can expand many states irrelevant to the goal;
          \item uses no estimate of remaining distance.
        \end{itemize}
      \end{block}
    \end{column}
  \end{columns}

  \begin{keyidea}
    Dijkstra is cost-aware, but it is still uninformed about where the goal lies.
  \end{keyidea}
\end{frame}

\sectionframe{Part 5: Changing the Cost Function Changes the Solution}

\begin{frame}[shrink=20]{Same Puzzle, Different Objective}
  Consider the sliding-tile puzzle from Lecture 5.

  \begin{columns}[T,onlytextwidth]
    \begin{column}{0.48\textwidth}
      \begin{block}{Uniform action model}
        Every legal tile move costs 1.

        \vspace{0.08in}
        Objective: minimize number of moves.
      \end{block}
    \end{column}

    \begin{column}{0.48\textwidth}
      \begin{block}{Weighted action model}
        Moving tile $k$ costs $k$.

        \vspace{0.08in}
        Objective: minimize sum of moved tile values.
      \end{block}
    \end{column}
  \end{columns}

  \begin{keyidea}
    States and legal transitions can remain identical while the preferred solution changes because the objective changed.
  \end{keyidea}
\end{frame}

\begin{frame}[shrink=20]{Micromouse: Cost Can Approximate Execution Time}
  A time-oriented action model might assign different costs to different motion primitives.

  \begin{center}
  \renewcommand{\arraystretch}{1.25}
  \begin{tabular}{l|c}
    \textbf{Motion primitive} & \textbf{Example modeled cost} \\\hline
    long straight segment & low cost per unit distance \\
    gentle curve & moderate cost \\
    sharp 90-degree turn & higher cost \\
    enter / exit diagonal & transition cost \\
    braking and re-acceleration & additional time cost
  \end{tabular}
  \end{center}

  \begin{keyidea}
    The search objective should reflect the performance metric we actually care about---for racing, that may be time rather than geometric distance.
  \end{keyidea}
\end{frame}

\begin{frame}[shrink=20]{State, Actions, Transitions, and Costs Work Together}
  \begin{columns}[T,onlytextwidth]
    \begin{column}{0.47\textwidth}
      \begin{block}{Simple maze model}
        \begin{itemize}
          \item State: cell
          \item Actions: orthogonal moves
          \item Transition: adjacent cell
          \item Cost: 1 per move
        \end{itemize}
      \end{block}
    \end{column}

    \begin{column}{0.47\textwidth}
      \begin{block}{Richer physical model}
        \begin{itemize}
          \item State: position, heading, velocity
          \item Actions: motion primitives
          \item Transition: kinematic/dynamic model
          \item Cost: predicted execution time
        \end{itemize}
      \end{block}
    \end{column}
  \end{columns}

  \begin{keyidea}
    Better planning often comes from improving the problem formulation, not merely replacing the search algorithm.
  \end{keyidea}
\end{frame}

\sectionframe{Part 6: Where Minimum-Cost Search Appears in AI}

\begin{frame}[shrink=20]{Robotics and Autonomous Navigation}
  \begin{columns}[T,onlytextwidth]
    \begin{column}{0.54\textwidth}
      Weighted graph search is a natural baseline when a robot must trade among:
      \begin{itemize}
        \item travel distance;
        \item estimated execution time;
        \item energy use;
        \item terrain difficulty;
        \item clearance or collision risk.
      \end{itemize}

      \begin{block}{Typical abstraction}
        roadmap / occupancy grid $\rightarrow$ weighted graph $\rightarrow$ minimum-cost route
      \end{block}
    \end{column}

    \begin{column}{0.40\textwidth}
    \end{column}
  \end{columns}
\end{frame}

\begin{frame}[shrink=20]{Autonomous Driving: Global Route Cost}
  A driving route may be evaluated using quantities such as:
  \begin{itemize}
    \item travel time;
    \item road length;
    \item speed limits;
    \item turn penalties;
    \item road class;
    \item predicted congestion or risk.
  \end{itemize}

  \begin{block}{Important distinction}
    Real autonomous-driving stacks use much richer models than plain Dijkstra. Dijkstra provides the foundational minimum-cost graph-search principle over a discretized route network.
  \end{block}

  \begin{keyidea}
    Weighted graph search often solves the global route problem; local motion planning must additionally respect vehicle dynamics and nearby obstacles.
  \end{keyidea}
\end{frame}

\begin{frame}[shrink=20]{Symbolic Planning with Nonuniform Actions}
  An AI planner may reason over actions such as:
  \begin{center}
    \texttt{drive} \qquad \texttt{fly} \qquad \texttt{charge} \qquad \texttt{purchase} \qquad \texttt{wait}
  \end{center}

  Different actions may incur different:
  \begin{itemize}
    \item durations;
    \item monetary costs;
    \item resource consumption;
    \item risk penalties.
  \end{itemize}

  The state-space graph can therefore be searched for a minimum-cost plan rather than a minimum-number-of-actions plan.

  \begin{keyidea}
    Uniform-cost search is a general planning primitive whenever action sequences have additive non-negative costs.
  \end{keyidea}
\end{frame}

\begin{frame}[shrink=20]{Game and Simulation AI}
  Navigation graphs and navigation meshes often encode nonuniform traversal preferences.

  \begin{center}
  \renewcommand{\arraystretch}{1.25}
  \begin{tabular}{l|c}
    \textbf{Region} & \textbf{Example cost} \\\hline
    open floor & 1 \\
    mud / difficult terrain & 3 \\
    exposed danger zone & 10
  \end{tabular}
  \end{center}

  An AI-controlled character may choose a geometrically longer route because it has lower modeled total cost.

  \begin{keyidea}
    ``Shortest'' is meaningful only after the system designer specifies the metric being minimized.
  \end{keyidea}
\end{frame}

\begin{frame}[shrink=20]{Knowledge Graphs: Useful Subroutine, Not the Whole System}
  Knowledge graphs represent:
  \begin{itemize}
    \item entities as vertices;
    \item relationships as edges;
    \item confidence, distance, or relation penalties as possible weights.
  \end{itemize}

  Weighted shortest-path computations can support tasks such as:
  \begin{itemize}
    \item finding low-cost relation chains;
    \item measuring graph proximity;
    \item generating candidate explanation paths;
    \item ranking graph-connected candidates.
  \end{itemize}

  \begin{alertblock}{Modern AI context}
    Entity linking and knowledge-graph reasoning are usually not ``just Dijkstra.'' Learned representations, retrieval models, language models, and graph neural methods often provide the dominant signal. Shortest-path search is better viewed as a reusable graph primitive inside a larger system.
  \end{alertblock}
\end{frame}

\begin{frame}[shrink=20]{Search Appears Inside Larger AI Architectures}
  \begin{center}
    \large perception / retrieval / learned model
    $\longrightarrow$
    weighted graph or planning problem
    $\longrightarrow$
    search
    $\longrightarrow$
    action or explanation
  \end{center}

  \begin{itemize}
    \item AI systems frequently combine learned components with classical algorithms.
    \item A learned model may estimate costs or construct the graph.
    \item A search algorithm may then enforce global consistency or find a minimum-cost sequence.
  \end{itemize}

  \begin{keyidea}
    Classical search and modern machine learning are complementary computational tools, not mutually exclusive definitions of AI.
  \end{keyidea}
\end{frame}

\sectionframe{Part 7: Why Dijkstra Is Still Uninformed}

\begin{frame}[shrink=20]{Cost-Aware Does Not Mean Goal-Informed}
  Dijkstra / UCS knows:
  \[
    g(n)=\text{cost already paid to reach }n.
  \]

  It does \textbf{not} know:
  \begin{itemize}
    \item how far $n$ is from the goal;
    \item whether $n$ points in a promising direction;
    \item whether a neighboring region is likely to contain a solution.
  \end{itemize}

  \begin{block}{Decision rule}
    \centering
    ``Which discovered state has been cheapest to reach so far?''
  \end{block}

  \begin{keyidea}
    Dijkstra uses path cost, but no estimate of remaining cost. It is therefore still an uninformed search strategy.
  \end{keyidea}
\end{frame}

\begin{frame}[shrink=20]{The Cost Contour Can Expand in the Wrong Direction}
  \begin{columns}[T,onlytextwidth]
    \begin{column}{0.54\textwidth}
      Suppose the goal lies to the east, but a large region to the south contains many very cheap edges.

      Dijkstra may expand much of that region first because:
      \[
        g(\text{southern states}) < g(\text{goal route}).
      \]

      Nothing in $g(n)$ says which direction leads toward the goal.
    \end{column}

    \begin{column}{0.40\textwidth}
    \end{column}
  \end{columns}

  \begin{keyidea}
    Optimality can require substantial exploration when the search has no estimate of remaining distance.
  \end{keyidea}
\end{frame}

\begin{frame}[shrink=20]{Lecture 5 vs. Lecture 6}
  \begin{center}
  \renewcommand{\arraystretch}{1.35}
  \begin{tabular}{l|c|c}
    & \textbf{Lecture 5} & \textbf{Lecture 6} \\\hline
    Action model & effectively equal steps & variable costs \\
    Objective & minimum depth & minimum path cost \\
    Main frontier & FIFO / stack / depth limit & priority queue \\
    Key strategies & BFS, DFS, IDDFS & Dijkstra / UCS \\
    Goal knowledge & none & none
  \end{tabular}
  \end{center}

  \begin{keyidea}
    We changed what ``best'' means, but we still have no domain-specific estimate of which frontier state is closest to success.
  \end{keyidea}
\end{frame}

\begin{frame}[shrink=20]{Search Quality Depends on the Model}
  The Micromouse story gives us a complete modeling ladder:

  \begin{enumerate}
    \item Change the \textbf{state}: orientation and velocity become relevant.
    \item Change the \textbf{actions}: diagonal and smooth motion becomes possible.
    \item Change the \textbf{transition model}: physical motion constraints matter.
    \item Change the \textbf{cost function}: fastest can replace shortest.
    \item Change the \textbf{frontier policy}: minimum-cost search becomes necessary.
  \end{enumerate}

  \begin{keyidea}
    An AI agent searches the world described by its model. Improving the model can change both which solutions exist and which solution counts as best.
  \end{keyidea}
\end{frame}

\begin{frame}[shrink=20]{Exit Ticket: Formulate the Right Search Problem}
  \begin{activity}
    An autonomous delivery robot can either cross an open courtyard or travel around it on a covered walkway. Rain makes the courtyard slippery and slows the robot.

    Specify:
    \begin{enumerate}
      \item what additional state or environment information may matter;
      \item how \texttt{ACTIONS(s)} could change;
      \item what the action cost should represent;
      \item why BFS may return a different route from Uniform-Cost Search.
    \end{enumerate}
  \end{activity}

  \begin{keyidea}
    Outcome: you can explain when a uniform-step search model is insufficient and formulate a minimum-cost search problem that reflects physical or task-specific constraints.
  \end{keyidea}
\end{frame}

\begin{frame}[shrink=20]{Next Time: Heuristic Search}
  Dijkstra asks:
  \begin{block}{Known cost}
    \centering
    How much has this path cost so far? \qquad $g(n)$
  \end{block}

  Suppose the agent could also estimate:
  \begin{block}{Remaining cost}
    \centering
    How much more might it cost to reach the goal? \qquad $h(n)$
  \end{block}

  Then search can use information about the goal rather than expanding solely by accumulated cost.

  \begin{keyidea}
    Next: Greedy Best-First Search, heuristic functions, and A$^*$.
  \end{keyidea}
\end{frame}

\end{document}
